% ---
% titre: Problème de proportionnalité - la provision de sucre
% theme: FA
% chapitre: Proportionnalité
% sous_chapitre: Résoudre des problèmes
% niveau: [10e]
% type: EVAL
% tags: [proportionnalite, p-question-TS]
% difficulte: 3
% corrige: true
% ---

\question[5]
Une réserve de sucre suffit à l'alimentation de 720 personnes pendant 24 jours. Que durerait la même provision s'il fallait nourrir 900 personnes ? Donne ta réponse en jour/heure/minute.

\vspace{10pt}{\footnotesize\raggedleft Espace pour ta démarche et tes calculs\par}
\begin{solutionorgrid}[180pt]

\vspace{5pt}
C'est une situation de proportionnalité inverse, plus il y a de personnes à nourrir, moins la réserve dure longtemps. \hfill \textbf{(1\,pt)}

\vspace{15pt}
Quantité totale de sucre en jours-personnes

\vspace{5pt}$720 \cdot 24  = 17 280$\,jours-personnes \hfill calcul \textbf{(0,5\,pt)} et réponse \textbf{(0,5\,pt)}

\vspace{15pt}
Durée pour 880 personnes

\vspace{5pt}$17\,280 : 900 = 19,2$\,jours \hfill calcul \textbf{(0,5\,pt)} et réponse \textbf{(0,5\,pt)}

\vspace{15pt}
Conversion en jours, heures, minutes

\vspace{5pt}Jours: $19$\,jours \textbf{(0,5\,pt)} \quad/\quad Heures: $0,2 \cdot 24 = 4,8$\,heures \textbf{(0,5\,pt)} \quad/\quad Minutes: $0,8 \cdot 60 = 48$\, minutes \textbf{(0,5\,pt)}

$\longrightarrow$ 19\,jours 4\,heures 48\,minutes
\end{solutionorgrid}

\begin{tcolorbox}[enhanced jigsaw, interior hidden, colframe=box, parbox=false]%
Réponse: 
\sol{\vspace{20pt}}{La réserve durerait 19 jours 4\,heures  et 48\,minutes pour nourrir 900 personnes. \textbf{(0,5pt)}}
\end{tcolorbox}



